% Formelsammlung zum Kompaktskript (zwei Seiten)
\documentclass[fontsize=10pt,a4paper,DIV=20,headsepline]{scrartcl}
% Gemeinsame Präambel für Skript.tex und die Einzelkapitel (einzelkapitel.tex)
% Gemeinsamer Stil der Übungsblätter und Tests (Kopie im Ordner stil/)
\makeatletter\def\input@path{{stil/}}\makeatother
\input{gdi-stil}
\input{header}
\usepackage{needspace}
\usepackage{multicol}
\usetikzlibrary{shapes.geometric, shapes.multipart, positioning}

%%%%%%%%%%%%%%%%%%%%%%%%%%%%%%%%%%%%%%%%%%%%%%%%%%
% Bausteine des Skripts
%%%%%%%%%%%%%%%%%%%%%%%%%%%%%%%%%%%%%%%%%%%%%%%%%%
% Merkkasten (wie "Kurz erklärt" auf den Arbeitsblättern)
\newtcolorbox{merke}[1][Merke]{enhanced, breakable, colback=diagramm!5, colframe=diagramm!80, boxrule=0.6pt, arc=1.5mm,
	left=3mm, right=3mm, top=1.5mm, bottom=1.5mm, title=#1, fonttitle=\bfseries\small, colbacktitle=diagramm!80, coltitle=white,
	before skip=3mm, after skip=3mm}
% Hinweis-Kasten in den Farben der Code-Kästen
\newtcolorbox{hinweis}{enhanced, sharp corners, boxrule=0pt, colback=codeHintergrund,
	borderline west={2pt}{0pt}{diagramm}, left=3mm, top=1.5mm, bottom=1.5mm, before skip=3mm, after skip=3mm}

% Die drei Teile jedes Kapitels
\newcommand{\teil}[1]{\needspace{9\baselineskip}\subsection*{\textcolor{codeSchluessel}{#1}}}
\newcommand{\worum}{\teil{Worum geht es?}}
\newcommand{\beispiele}{\teil{Beispiele}}
\newcommand{\aufgaben}{\teil{Aufgaben}}
\newcommand{\ueben}{\teil{Zum Weiterüben}}

% Beispiel mit Titel
\newcounter{beispiel}[section]
\newcommand{\beispiel}[1]{\par\medskip\needspace{6\baselineskip}\refstepcounter{beispiel}%
	\noindent\textbf{\textcolor{diagramm}{Beispiel \thesection.\arabic{beispiel}:} #1}\par\nopagebreak\smallskip}

% Aufgabe mit Titel; optionale Markierung: \K = wie in der Klausur, \S = anspruchsvoll
\newcounter{aufgabe}[section]
\renewcommand{\theaufgabe}{\thesection.\arabic{aufgabe}}
\newcommand{\K}{\ \tcbox[on line, sharp corners, boxrule=0pt, colback=diagramm!80, coltext=white, left=1mm, right=1mm, top=0.2mm, bottom=0.2mm, fontupper=\sffamily\bfseries\scriptsize]{KLAUSUR}}
\renewcommand{\S}{\ \textcolor{diagramm}{$\bigstar$}}
\newcommand{\aufgabe}[2][]{\par\medskip\needspace{4\baselineskip}\refstepcounter{aufgabe}%
	\noindent\textbf{Aufgabe \theaufgabe: #2}#1\par\nopagebreak\smallskip}

% Tabellen
\newcolumntype{C}[1]{>{\centering\arraybackslash}p{#1}}
\newcolumntype{L}[1]{>{\raggedright\arraybackslash}p{#1}}
\newcommand{\leer}{\rule[-0.6em]{0pt}{2.0em}}
\newcommand{\luecke}[1][3cm]{\rule[-0.4ex]{#1}{0.4pt}}
\newcommand{\code}[1]{\texttt{#1}}
% In Schreibmaschinenschrift keine Ligaturen (-- bleibt --) und " bleibt ein Anführungszeichen
\DisableLigatures{encoding = *, family = tt*}
\AtBeginDocument{\shorthandoff{"}}

% Flussdiagramme (wie in den Tests zu Blatt 4)
\tikzset{
	startstop/.style={rectangle, rounded corners=3mm, draw=diagramm, semithick, fill=diagramm!15, minimum width=2cm, minimum height=7mm, font=\small},
	io/.style={trapezium, trapezium left angle=75, trapezium right angle=105, draw=diagramm, semithick, fill=white, minimum height=7mm, font=\small, align=center},
	process/.style={rectangle, draw=diagramm, semithick, fill=white, minimum width=2.6cm, minimum height=7mm, font=\small, align=center},
	decision/.style={diamond, aspect=2.4, draw=diagramm, semithick, fill=diagramm!8, inner sep=1pt, font=\small, align=center},
	arrow/.style={-{Latex[length=2mm]}, draw=diagramm, thick},
	lbl/.style={font=\footnotesize, fill=white, inner sep=1.5pt},
	% UML-Klassen
	klasse/.style={rectangle split, rectangle split parts=3, draw=diagramm, semithick, fill=diagramm!5, font=\ttfamily\footnotesize, align=left, rectangle split part align={center,left,left}, minimum width=3.2cm},
	erbt/.style={-{Triangle[open, length=3mm, width=3mm]}, draw=diagramm, semithick},
	assoz/.style={draw=diagramm, semithick}
}

% Kopf- und Fußzeile
\clearpairofpagestyles
\ihead{GDI -- Kompaktskript}
\ohead{\rightmark}
\cfoot{\pagemark}
\automark{section}
\setkomafont{pageheadfoot}{\normalfont\small\color{codeKommentar}}
\setcounter{tocdepth}{1}
\emergencystretch=3em

\ihead{GDI -- Formelsammlung}
\ohead{Grundlagen der Informatik}
\setlength{\columnsep}{5mm}
\setlength{\parskip}{0pt}

% Kasten für einen Themenblock
\newtcolorbox{fs}[1]{enhanced, colback=diagramm!4, colframe=diagramm!70, boxrule=0.5pt, arc=1mm,
	left=1.5mm, right=1.5mm, top=0.8mm, bottom=0.8mm, title={#1}, fonttitle=\bfseries\footnotesize,
	colbacktitle=diagramm!70, coltitle=white, toptitle=0.2mm, bottomtitle=0.2mm,
	before skip=1.2mm, after skip=1.2mm, fontupper=\footnotesize}
% kompakte Liste
\newenvironment{kl}{\begin{itemize}[leftmargin=3mm, itemsep=0.1mm, topsep=0.3mm, parsep=0pt, label=\textcolor{diagramm}{\textbullet}]}{\end{itemize}}
% kurzer Code-Block ohne Zeilennummern
\lstdefinestyle{fsCode}{style=gdiJava, numbers=none, basicstyle=\ttfamily\scriptsize, aboveskip=0.5mm, belowskip=0mm, xleftmargin=0mm}
\lstnewenvironment{fscode}{\lstset{style=fsCode}}{}

\begin{document}
\begin{multicols}{2}
\raggedcolumns

\begin{fs}{Bits und Bytes}
\begin{kl}
\item $n$ Bit ergeben $2^n$ Zustände. Für $k$ Zustände: kleinstes $n$ mit $2^n \geq k$.
\item 1 Byte $=$ 8 Bit, \ 1 Wort $=$ 2 Byte, \ 1 Doppelwort $=$ 4 Byte
\item 1 KB $= 2^{10}$ B, \ 1 MB $= 2^{20}$ B, \ 1 GB $= 2^{30}$ B, \ 1 TB $= 2^{40}$ B
\end{kl}
\smallskip
{\scriptsize\setlength{\tabcolsep}{2.2pt}
\begin{tabular}{@{}l|*{9}{r}@{}}
$n$ & 0 & 1 & 2 & 3 & 4 & 5 & 6 & 7 & 8 \\
$2^n$ & 1 & 2 & 4 & 8 & 16 & 32 & 64 & 128 & 256 \\ \hline
$n$ & 9 & 10 & 11 & 12 & 13 & 14 & 15 & 16 & 20 \\
$2^n$ & 512 & 1024 & 2048 & 4096 & 8192 & 16384 & 32768 & 65536 & $\approx 10^6$
\end{tabular}}
\end{fs}

\begin{fs}{Zahlensysteme}
\begin{kl}
\item $(z_k \ldots z_1 z_0)_b = z_k \cdot b^k + \ldots + z_1 \cdot b + z_0$, \ Ziffern $0$ bis $b - 1$
\item \textbf{Divisionsmethode:} wiederholt durch $b$ teilen, Reste \textbf{von unten nach oben} lesen
\item \textbf{Subtraktionsmethode:} Stellenwerte von groß nach klein abziehen (passt: 1, sonst 0), Ziffern \textbf{von oben nach unten} lesen
\item Dual $\leftrightarrow$ oktal: Dreiergruppen, \ dual $\leftrightarrow$ hex: Vierergruppen, jeweils \textbf{von rechts}
\end{kl}
\smallskip
{\scriptsize\setlength{\tabcolsep}{2.4pt}
\begin{tabular}{@{}l|*{8}{c}@{}}
hex & 0 & 1 & 2 & 3 & 4 & 5 & 6 & 7 \\
dual & 0000 & 0001 & 0010 & 0011 & 0100 & 0101 & 0110 & 0111 \\ \hline
hex & 8 & 9 & A & B & C & D & E & F \\
dual & 1000 & 1001 & 1010 & 1011 & 1100 & 1101 & 1110 & 1111
\end{tabular}}
\end{fs}

\begin{fs}{Zweierkomplement mit $n$ Bit}
\begin{kl}
\item \textbf{Negieren:} alle Bits umdrehen, dann 1 addieren. \ Rücklesen: ebenso, Minus davor.
\item Vorderstes Bit: 0 positiv, 1 negativ. Es gibt nur eine Null.
\item Wertebereich $-2^{n-1}$ bis $2^{n-1} - 1$ \ (8 Bit: $-128 \ldots 127$, \ 16 Bit: $-32\,768 \ldots 32\,767$)
\item $a - b = a + (-b)$. Ein Übertrag über die feste Breite fällt weg.
\item \textbf{Überlauf}, wenn Übertrag \emph{in} die Vorzeichenstelle $\neq$ Übertrag \emph{aus} ihr.
\end{kl}
\end{fs}

\begin{fs}{Nachkommastellen und Gleitkomma}
\begin{kl}
\item $2^{-1} = 0{,}5$, \ $2^{-2} = 0{,}25$, \ $2^{-3} = 0{,}125$, \ $2^{-4} = 0{,}0625$, \ $2^{-5} = 0{,}03125$
\item Dezimal $\to$ dual: Nachkommateil immer wieder $\cdot\,2$, Ziffer vor dem Komma notieren, \textbf{von oben nach unten} lesen
\item Normieren auf $1{,}M \cdot 2^E$. Gespeichert: Vorzeichen, $E + \mathit{Bias}$, $M$ ohne führende 1
\item \code{float}: $1 + 8 + 23$ Bit, Bias 127 \quad \code{double}: $1 + 11 + 52$ Bit, Bias 1023
\item $0{,}1$, $0{,}2$, \ldots\ sind dual periodisch $\Rightarrow$ Gleitkommazahlen nie mit \code{==} vergleichen
\end{kl}
\end{fs}

\begin{fs}{Textcodierung}
\begin{kl}
\item \textbf{ASCII} (7 Bit): \code{A} $= 41_{16} = 65$, \ \code{a} $= 61_{16} = 97$, \ \code{0} $= 30_{16} = 48$; \ klein $=$ groß $+\ 20_{16}$
\item \textbf{ISO 8859-1} (8 Bit): \code{Ä} C4, \code{Ö} D6, \code{Ü} DC, \code{ä} E4, \code{ö} F6, \code{ü} FC, \code{ß} DF
\item \textbf{UTF-8} (1 bis 4 Byte): \ \code{0xxxxxxx} \ | \ \code{110xxxxx 10xxxxxx} \ | \ \code{1110xxxx 10xxxxxx 10xxxxxx}
\item \textbf{UTF-16}: 2 (selten 4) Byte. In Java ist ein \code{char} 2 Byte groß.
\end{kl}
\end{fs}

\begin{fs}{Boolesche Algebra}
\begin{center}{\scriptsize\setlength{\tabcolsep}{4pt}
\begin{tabular}{@{}cc|cccc@{}}
$a$ & $b$ & $a \wedge b$ & $a \vee b$ & $a \oplus b$ & $\neg a$ \\ \hline
0 & 0 & 0 & 0 & 0 & 1 \\
0 & 1 & 0 & 1 & 1 & 1 \\
1 & 0 & 0 & 1 & 1 & 0 \\
1 & 1 & 1 & 1 & 0 & 0
\end{tabular}}\end{center}\vspace{-2mm}

\begin{kl}
\item De Morgan: $\neg(a \wedge b) = \neg a \vee \neg b$, \ $\neg(a \vee b) = \neg a \wedge \neg b$
\item Absorption: $a \vee (a \wedge b) = a$, \ $a \wedge (a \vee b) = a$
\item Distributiv: $a \wedge (b \vee c) = (a \wedge b) \vee (a \wedge c)$, \ $a \vee (b \wedge c) = (a \vee b) \wedge (a \vee c)$
\item $a \wedge \neg a = 0$, \ $a \vee \neg a = 1$, \ $a \wedge 1 = a$, \ $a \vee 0 = a$, \ $a \wedge 0 = 0$, \ $a \vee 1 = 1$, \ $\neg\neg a = a$
\item Eine Wahrheitstafel mit $n$ Variablen hat $2^n$ Zeilen.
\end{kl}
\end{fs}

\begin{fs}{Rechner und Algorithmen}
\begin{kl}
\item \textbf{Von Neumann:} Rechenwerk (ALU), Steuerwerk, Speicherwerk, Ein-/Ausgabewerk, verbunden durch den Bus. Programme \emph{und} Daten im selben Speicher.
\item Schichten: Hardware -- Firmware -- Betriebssystem -- Oberfläche -- Anwendungen
\item \textbf{terminierend:} endet nach endlich vielen Schritten. \ \textbf{determiniert:} gleiche Eingabe, gleiches Ergebnis. \ \textbf{deterministisch:} nächster Schritt stets eindeutig. \ \textbf{total korrekt} $=$ partiell korrekt $+$ terminierend
\item \textbf{Compiler:} übersetzt vor der Ausführung. \ \textbf{Interpreter:} zur Laufzeit. \ Java: \code{javac} $\to$ Byte-Code $\to$ JVM
\end{kl}
\end{fs}

\begin{fs}{Grammatiken und EBNF}
\begin{kl}
\item $G = (N, T, P, S)$: Nichtterminale, Terminale, Produktionen, Startsymbol
\item EBNF: \code{"a"} Terminal, \ \code{A B} Folge, \ \code{A | B} oder, \ \code{[A]} 0- oder 1-mal, \ \code{\{A\}} 0-mal oder öfter, \ \code{(A)} Gruppe
\item $a^n b^n$, $n \in \mathbb{N}$: $\{(S, aSb), (S, ab)\}$; \ mit $n \in \mathbb{N}_0$: $\{(S, aSb), (S, \varepsilon)\}$
\end{kl}
\smallskip
{\scriptsize\setlength{\tabcolsep}{3pt}
\begin{tabular}{@{}p{2.3cm}p{5.6cm}@{}}
\textbf{Typ 3} regulär & rechts höchstens ein Nichtterminal, immer am Rand; \emph{endlicher Automat} \\
\textbf{Typ 2} kontextfrei & links genau ein Nichtterminal; \emph{Kellerautomat} \\
\textbf{Typ 1} kontextsensitiv & Regeln verkürzen nicht: $|p| \leq |q|$ \\
\textbf{Typ 0} allgemein & beliebige Regeln; \emph{Turingmaschine}
\end{tabular}}
\end{fs}

\begin{fs}{Theoretische Informatik}
\begin{kl}
\item Endlicher Automat $M = (S, A, \delta, s_0, F)$: akzeptiert, wenn er nach dem letzten Zeichen in $F$ ist.
\item Reguläre Ausdrücke: $e + f$ (oder), $ef$ (Folge), $e^*$ (beliebig oft, auch nie); Bindung $^*$ vor $\cdot$ vor $+$
\item These von Church; Halteproblem und Satz von Rice: nicht entscheidbar
\item $\{\mathit{VOR}\}\ \mathit{ANW}\ \{\mathit{NACH}\}$: Gilt VOR und terminiert ANW, gilt danach NACH.
\end{kl}
\end{fs}

\end{multicols}
\newpage
\begin{multicols}{2}
\raggedcolumns

\begin{fs}{Java: Datentypen}
{\scriptsize\setlength{\tabcolsep}{3pt}
\begin{tabular}{@{}ll@{\hspace{3mm}}ll@{}}
\code{byte} & 1 Byte, $-128 \ldots 127$ & \code{short} & 2 Byte, $-32\,768 \ldots 32\,767$ \\
\code{int} & 4 Byte, ca. $\pm 2{,}1 \cdot 10^9$ & \code{long} & 8 Byte, ca. $\pm 9{,}2 \cdot 10^{18}$ \\
\code{float} & 4 Byte & \code{double} & 8 Byte \\
\code{char} & 2 Byte, Unicode & \code{boolean} & \code{true} oder \code{false}
\end{tabular}}
\begin{kl}
\item Implizit: \code{byte} $\to$ \code{short} $\to$ \code{int} $\to$ \code{long} $\to$ \code{float} $\to$ \code{double}, \ \code{char} $\to$ \code{int}. \ Andersherum nur mit Cast: \code{(int) 3.9} ist 3.
\item Ganzzahl durch Ganzzahl ist ganzzahlig: \code{7 / 2} ist 3, \ \code{7 \% 2} ist 1, \ \code{7 / 2.0} ist 3.5
\item Literale: \code{42} int, \code{42L} long, \code{1.5} double, \code{1.5f} float, \code{'x'} char, \code{"x"} String
\item Lokale Variablen vor dem ersten Lesen initialisieren. \ Konstante: \code{final double PI = 3.14;}
\end{kl}
\end{fs}

\begin{fs}{Java: Operatoren}
\begin{kl}
\item Priorität: \code{++ -- !} Cast \ > \ \code{* / \%} \ > \ \code{+ -} \ > \ \code{< <= > >=} \ > \ \code{== !=} \ > \ \code{\&\&} \ > \ \code{||} \ > \ \code{?:} \ > \ \code{= +=} \ldots
\item \code{x++}: alten Wert verwenden, dann erhöhen. \ \code{++x}: erst erhöhen, dann verwenden.
\item \code{a op= b} bedeutet \code{a = a op b}; der alte Wert von \code{a} wird zuerst gemerkt.
\item \code{\&\&}, \code{||}: rechte Seite nur, wenn nötig. \ \code{\&}, \code{|}: immer beide Seiten.
\item Von links: \code{"x" + 1 + 2} ist \code{"x12"}, \ \code{1 + 2 + "x"} ist \code{"3x"}
\end{kl}
\end{fs}

\begin{fs}{Java: Ein- und Ausgabe, Kontrollstrukturen}
\begin{fscode}
import java.util.Scanner;
Scanner sc = new Scanner(System.in);
int n = sc.nextInt();   double d = sc.nextDouble();
String s = sc.next();
System.out.println("n = " + n);    // print: ohne Umbruch
double r = Math.round(x * 100) / 100.0;   // 2 Stellen

if (a > 0) { ... } else if (a < 0) { ... } else { ... }
switch (n) {
	case 1: ...; break;
	case 2: case 3: ...; break;
	default: ...
}
while (b) { ... }        // evtl. nie
do { ... } while (b);    // mindestens einmal
for (int i = 0; i < n; i++) { ... }
for (int w : feld) { ... }  // nur lesen
\end{fscode}
\end{fs}

\begin{fs}{Java: Arrays und Methoden}
\begin{fscode}
int[] f = new int[5];   int[] g = {3, 1, 4};
int[][] m = new int[2][3];
static int max(int a, int b) { return a > b ? a : b; }
\end{fscode}
\begin{kl}
\item Index \code{0} bis \code{f.length - 1}, sonst \code{ArrayIndexOutOfBoundsException}. Vorbelegung 0, 0.0, \code{false}, \code{null}.
\item \code{m.length} ist 2 (Zeilen), \code{m[0].length} ist 3. \ \code{b = a} kopiert nur die Referenz, \code{a == b} vergleicht Referenzen.
\item Signatur $=$ Name $+$ Parametertypen (ohne Rückgabetyp). Überladen: gleicher Name, andere Parameter.
\item Call-by-value: Parameter sind Kopien; bei Arrays und Objekten wird die Referenz kopiert.
\end{kl}
\end{fs}

\begin{fs}{Rekursion und Komplexität}
\begin{kl}
\item Rekursion braucht einen \textbf{Basisfall} und einen \textbf{Rekursionsschritt} zu einem kleineren Problem. Auswerten: Aufrufe notieren, Ergebnisse von innen nach außen einsetzen.
\item Endrekursiv: Der Aufruf ist die letzte Aktion, z.\,B. \code{return fak(n - 1, akk * n);}
\item $O(1) \subset O(\log n) \subset O(n) \subset O(n \log n) \subset O(n^2) \subset O(n^3) \subset O(2^n)$
\item Konstanten und kleinere Summanden entfallen: $O(3n^2 + 5n) = O(n^2)$
\item Schleife: Durchläufe $\times$ Rumpf. Geschachtelt: multiplizieren. Nacheinander: der größte zählt. Halbieren pro Schritt: $\log n$.
\end{kl}
\end{fs}

\begin{fs}{Objektorientierung}
\begin{fscode}
public class Konto {
	private static int naechsteNr = 1;   // für eindeutige Nummern
	private int nr;
	private double saldo;
	public Konto(double saldo) {   // Konstruktor
		this.nr = naechsteNr++;
		this.saldo = saldo;
	}
	public double getSaldo() { return saldo; }
}
public abstract class Figur {
	public abstract double getFlaeche();
}
public class Kreis extends Figur implements Skalierbar {
	public Kreis(double r) { super(); ... }   // zuerst!
	@Override public double getFlaeche() { ... }
}
\end{fscode}
\begin{kl}
\item UML: \code{-} private, \code{+} public, \code{\#} protected, unterstrichen $=$ \code{static}, kursiv bzw. \{abstract\} $=$ abstrakt. Hohle Pfeilspitze zeigt auf die Oberklasse. Multiplizität \code{1}, \code{*}, \code{0..n}.
\item Unveränderlich: im Konstruktor setzen, keine \code{set}-Methode. Später festgelegt: \code{set}-Methode.
\item Abstrakte Klasse: keine Objekte. Interface: nur Methoden ohne Rumpf und Konstanten; mehrere pro Klasse.
\item \code{Figur f = new Kreis(1);} (Up-Cast) \ \code{Kreis k = (Kreis) f;} (Down-Cast) \ \code{f instanceof Kreis}
\end{kl}
\end{fs}

\begin{fs}{Strings und Wrapper}
\begin{kl}
\item \code{length()}, \code{charAt(i)}, \code{indexOf(t)} ($-1$, falls nicht enthalten), \code{substring(von, bis)} (\code{bis} exklusiv), \code{equals(t)}, \code{replace(a, b)}, \code{endsWith(t)}, \code{toCharArray()}
\item Strings sind unveränderlich: \code{s = s.replace('a', 'b');} \ Vergleich nie mit \code{==}, immer mit \code{equals}.
\item \code{Integer.parseInt("42")}, \ \code{Integer.toString(42)}, \ \code{Integer.MAX\_VALUE}, \ \code{Integer i = 5;} (Autoboxing)
\end{kl}
\end{fs}

\end{multicols}
\end{document}
